\documentclass{beamer}

\usepackage{beamerthemesplit}
\usetheme{Singapore} %Copenhagen}
%\usecolortheme{whale}

%\usepackage[T2A]{fontenc}
%\usepackage[utf8]{inputenc}
%\usepackage[russian]{babel}

\usepackage[main=russian,english]{babel}   %% загружает пакет многоязыковой вёрстки
\usepackage{fontspec}      %% подготавливает загрузку шрифтов Open Type, True Type и др.
\defaultfontfeatures{Ligatures={TeX},Renderer=Basic}  %% свойства шрифтов по умолчанию
\setmainfont{Times New Roman} %% задаёт основной шрифт документа
%\usefonttheme{professionalfonts}% SOLUTION
\usefonttheme{serif}

\usepackage{hyperref}
\usepackage{textcomp}
\usepackage{amssymb,amsmath}
%\usepackage{animate}
%\usepackage{longtable}
\usepackage{xcolor}

%\usepackage{pgffor}
\usepackage{enumitem}
\usepackage[export]{adjustbox}

\newcounter{N}

%% Форматирование окружения itemize
%\usepackage{ragged2e}
%\let\olditem\item
%\renewcommand\item{\olditem\justifying}

\usepackage{ mathrsfs }
\newcommand{\Rho}{\mathscr{P}}

\renewcommand{\Re}{\operatorname{Re}}

\newcommand{\argxi}{(\xi^1,\xi^2,\xi^3)}
\newcommand{\argx}{(x^1,x^2,x^3)}

\newcommand{\argxiv}{(\vec{\xi})}
\newcommand{\argxv}{(\vec{x})}


\newcommand{\argxbarn}{(\bar{x}^1,\bar{x}^2,\ldots, \bar{x}^n)}
\newcommand{\argxn}{(x^1, x^2,\ldots, x^n)}

\newcommand{\argtxi}{(t, \xi^1,\xi^2,\xi^3)}
\newcommand{\argtoxi}{(t_0, \xi^1,\xi^2,\xi^3)}

\newcommand{\argtxiv}{(t, \vec{\xi})}
\newcommand{\argtoxiv}{(t_0, \vec{\xi})}


\newcommand{\argtx}{(t, x^1,x^2,x^3)}
\newcommand{\argtox}{(t_0, x^1,x^2,x^3)}

\newcommand{\argtxv}{(t, \vec{x})}
\newcommand{\argtoxv}{(t_0, \vec{x})}


\newcommand{\pd}[2]{\frac{\partial #1}{\partial #2}}
\newcommand{\pdk}[2]{\frac{\partial^2 #1}{\partial #2^2}}

\newcommand{\od}[2]{\frac{d #1}{d #2}}

\newcommand{\grad}{\operatorname{grad}}
\newcommand{\rot}{\operatorname{rot}}
\newcommand{\divo}{\operatorname{div}}

\title[]{Трёхмерные осесимметричные потенциальные течения идеальной жидкости}

\author[]{ {\em Верещагин Антон Сергеевич}
\\
канд. физ.-мат. наук, старший преподаватель\\
\bigskip
Кафедра аэрофизики и газовой динамики ФФ НГУ}

\usebackgroundtemplate{\includegraphics[width=\paperwidth]{../img/background.png}}

\begin{document}
	
\frame{\titlepage}


\frame{
	\frametitle{Аннотация}
	\parbox{\textwidth}{
		
	}
}

\frame{
	\frametitle{ Основные определения}
	
	\begin{exampleblock}{Определение}
	\parbox{\textwidth}{
		Течение называется \alert{осесимметричным}, если существует такая прямая $l$, что во всех плоскостях, проходящих через $l$ картина течения одинакова и траектория жидкой частицы лежит в полуплоскостях, проходящих через $l$.
		
	}
	\end{exampleblock}

	\begin{exampleblock}{Определение}
		\parbox{\textwidth}{
			Течение называется \alert{потенциальным}, если в некоторой области пространства можно определить потенциал $\varphi(t,x,y,z)$, такой что
			\[
			\vec{v} = \grad \varphi.
			\]
			
		}
	\end{exampleblock}

}



\frame{
	\frametitle{Основные уравнения}
	
	\parbox{\textwidth}{
	Для трёхмерных потенциальных течений идеальной жидкости определённых в некоторой области пространства справедливы следующие уравнения.	
	}
	
	\bigskip
	
	
	\begin{exampleblock}{Уравнение неразрывности}
		\parbox{\textwidth}{
			\[
				\divo\vec{v}=\Delta\varphi = 0,\quad
				\vec{v}=\nabla \varphi
			\]
		}
	\end{exampleblock}
	
	\begin{exampleblock}{Интеграл Коши}
		\parbox{\textwidth}{
			\[
			\frac{\nabla\varphi^2}{2} + \frac{p}{\rho} = f(t)\footnote{Считаем, что поле внешних сил отстутствует}
			\]
			
		Интеграл Коши позволяет найти распределение давления по заданному потенциалу, определённому из уравнения неразрывности ($\rho=const$).
		}
	\end{exampleblock}
}

\frame{
	\frametitle{ Уравнение неразрывности в различных системах координат}
	
	\begin{exampleblock}{Сферическая система координат $r$, $\theta$, $\lambda$}
		\parbox{\textwidth}{
			\[
			\frac{1}{r^2 \sin\theta}
			\left\{
			\pd{}{r}\left( r^2 \sin\theta \pd{\varphi}{r} \right) + 
			\pd{}{\theta} \left( \sin\theta \pd{\varphi}{\theta} \right) +
			\pd{}{\lambda}\left( \frac{1}{\sin\theta} \pd{\varphi}{\lambda} \right)
			\right\}
			 = 0,
			\]
			где
			\[
			v_r = \pd{\varphi}{r},\quad
			v_\theta = \frac{1}{r}\pd{\varphi}{\theta},\quad
			v_\lambda = \frac{1}{r\sin\theta} \pd{\varphi}{\lambda}.
			\]
		}
	\end{exampleblock}
	
	\begin{exampleblock}{}
		\parbox{\textwidth}{
			В случае осесимметричного течения можно пренебречь зависимостью $\varphi$ от $\lambda$.
		}
	\end{exampleblock}
}

\frame{
	\frametitle{ Уравнение неразрывности в различных системах координат}
	
	\begin{exampleblock}{Цилиндрическая система координат $r$, $\theta$, $z$}
		\parbox{\textwidth}{
			\[
			\frac{1}{r}
			\left\{
			\pd{}{r}\left( r v_r \right) + 
			\pd{}{\theta} v_\theta +
			\pd{}{\lambda}\left( r v_z \right)
			\right\}
			= 0,
			\]
			где
			\[
			v_r = \pd{\varphi}{r},\quad
			v_\theta = \frac{1}{r}\pd{\varphi}{\theta},\quad
			v_z =  \pd{\varphi}{z}.
			\]
		}
	\end{exampleblock}
	
	\begin{exampleblock}{}
		\parbox{\textwidth}{
			В случае осесимметричного течения можно пренебречь зависимостью $\varphi$ от $\theta$.
		}
	\end{exampleblock}
}


\end{document}